\begin{figure}[tbp]
\centering
\begin{tikzpicture}[x=0.9cm,y=0.9cm,>=Stealth,
 every node/.style={font=\small}]
 \definecolor{flowblue}{RGB}{25,108,105}
 \foreach \a/\b/\lab in {0/1.3/I_1,2.5/3.8/I_j,3.8/5.1/I_{j+1},6.4/7.7/I_n}
  {\draw[fill=gray!7] (\a,0) rectangle (\b,0.75);
   \node at ({(\a+\b)/2},0.375) {$\lab$};}
 \node at (1.9,0.375) {$\cdots$};
 \node at (5.75,0.375) {$\cdots$};
 \draw[fill=flowblue!10] (7.7,0) rectangle (9.5,0.75);
 \node at (8.6,0.375) {$I_1+\beta$};
 \node at (10.05,0.375) {$\cdots$};
 \draw[->] (-0.25,-0.25)--(10.65,-0.25) node[right] {$t$};
 \draw (0,-0.2)--(0,-0.3) node[below] {$0$};
 \draw (7.7,-0.2)--(7.7,-0.3) node[below] {$\beta$};
 \draw[dashed] (7.7,-0.05)--(7.7,2.1);
 \draw[->,thick,flowblue] (3.15,0.9) to[out=70,in=110] (4.45,0.9);
 \node[align=center,text=flowblue] at (3.8,1.9) {propagate, then average\\on regular slices};
 \draw[->,thick,flowblue] (7.05,0.9) to[out=70,in=110] (8.6,0.9);
 \node[align=center,text=flowblue] at (8.5,2.55) {continue through\\the next period};
 \node[align=center] at (5.1,-1.3)
  {$\langle m,y\rangle=\langle m,y-U_d^*y\rangle+\langle U_dm,y\rangle$};
\end{tikzpicture}
\caption{Averaging the source mass from one block in the next block.
The last block uses the first block of a copied period. The displayed identity
separates the pairing of its terminal mass $m$ into a term controlled by contraction energy
and a term whose propagation can be split at the copied averaging slices.
The averaging windows are disjoint modulo one period, so their Dirichlet
energies are summed with a constant total cost. Block widths are schematic.}
\label{fig:period-averaging}
\end{figure}
